% Agent 24: Pages 432--437 (PDF pp.46--48)
% Sections 5.7 (cont.) Equations of Vertical Structure -- boundary conditions
%          5.8 Approximate Version of Vertical Structure
%          5.9 Explicit Models for Disk (outer region and middle region)
% Equations (5.7.1)--(5.9.9)
% =============================================================================

%% ----------------------------------------------------------------
%% 5.7  Equations of Vertical Structure (continuation from p.~431)
%% ----------------------------------------------------------------

% [Page 432 -- continuation of Section 5.7]

of special relativity (equivalence principle).  And because $\Pi \sim Tk/\rho_0 \ll 1$ those
special relativistic laws take on their standard Newtonian forms.  At least they
do so if one uses the correct Kerr-metric value for the tidal gravitational
acceleration which compresses the disk into its ``pancake'' shape:

(``acceleration of gravity'') $\equiv g = R^{\hat{z}}{}_{\hat{0}\hat{z}\hat{0}}\,z$
%
\begin{equation}
\tag{5.7.1}
\end{equation}

(cf.\ eq.~5.2.24).  An explicit expression for $R^{\hat{z}}{}_{\hat{0}\hat{z}\hat{0}}$
can be obtained by a transformation of the ``LNRF'' components,
$R^{(\hat{\alpha})}{}_{(\hat{\beta})(\hat{\gamma})(\hat{\delta})}$, of the Riemann tensor as
given in Bardeen, Press and Teukolsky~(1972)~\cite{BardeenPressTeukolsky1972}:
%
% NOTE: Eq.~(5.7.2) involves a complex expression for R_{\hat{0}\hat{z}\hat{0}\hat{z}}
% in terms of Kerr metric functions from Bardeen, Press \& Teukolsky (1972).
% The exact combination of script correction factors is difficult to read
% from the scan; the leading terms are:
\begin{equation}
\tag{5.7.2}
R_{\hat{0}\hat{z}\hat{0}\hat{z}}
= \frac{M}{r^3}
  \left[
    \frac{(r^2 + a^2)^2 + 2\Delta\, a^2}{(r^2 + a^2)^2 - \Delta\, a^2}
  \right]
  \frac{\mathscr{A}^2\,\mathscr{D}}
       {\mathscr{E}\,\mathscr{F}}
\end{equation}

Thus, a Newtonian astrophysicist can build a relativistically correct vertical
structure without knowing any relativity at all.  He need merely follow standard
Newtonian procedures and theory; but he must use expression~(5.7.1) for the
vertical acceleration of gravity, rather than the Newtonian formula $g = (M/r^3)\,z$:

We shall describe the vertical structure in terms of the following functions of
height---which we tacitly assume have all been averaged over $t$, $r$, $\varphi$ by the
methods of equation~(5.5.2):
%
\begin{alignat}{2}
&\rho(r,z) &&= \text{density of rest mass;} \notag \\
&p(r,z)   &&= \text{vertical pressure, } T_{\hat{z}\hat{z}}; \notag \\
&t_{\varphi r}(r,z) &&= \text{shear stress, } T_{\hat{z}\hat{z}}; \tag{5.7.3} \\
&T(r,z)   &&= \text{temperature;} \notag \\
&q^z(r,z) &&= \text{flux of energy;} \notag \\
&K(r,z)   &&= \text{Rosseland mean opacity.} \notag
\end{alignat}

These 6 vertical structure functions are governed by the following 6 ``equations
of vertical structure'':

\paragraph{Vertical pressure balance}
%
\begin{equation}
\tag{5.7.4a}
\frac{dp}{dz}
= \rho_0 \times (\text{``acceleration of gravity''})
= \rho_0\, R^{\hat{z}}{}_{\hat{0}\hat{z}\hat{0}}\, z
\end{equation}

\paragraph{Sources of viscosity}
%
\begin{equation}
\tag{5.7.4b}
t_{\varphi r}
= \left\{
    \begin{tabular}{@{}l@{}}
    some explicit expression which depends on\\
    the explicit assumptions of the model
    \end{tabular}
  \right\}
\end{equation}

\paragraph{Energy generation}
[cf.\ eq.~(5.6.11)]
%
\begin{equation}
\tag{5.7.4c}
\frac{dq^z}{dz}
= -2\sigma^{(\mathrm{SG})}_{\varphi r}\, t_{\varphi r}
= \tfrac{3}{2}(M/r^3)^{1/2}\, t_{\varphi r}\,
  \mathscr{G}^{-1}\,\mathscr{Q}.
\end{equation}

\paragraph{Energy transport}
%
\begin{equation}
\tag{5.7.4d}
q^z
= \left\{
    \begin{tabular}{@{}l@{}}
    some explicit expression which depends on the\\
    nature of the transport assumed by the model
    \end{tabular}
  \right\}.
\end{equation}
%
The energy transport may be by radiative diffusion in some models, but by
turbulent gas motions in others.  All models to date have assumed radiative
transfer with large optical depth, $\tau_{es} + \tau \gg 1$ (see~\S\ref{sec:2.6}).  In this case
%
\begin{equation}
\tag{5.7.4d$'$}
q^z = -\frac{1}{K\rho_0}\,\frac{d}{dz}\bigl(\tfrac{1}{3}b\,T^4\bigr).
\end{equation}

\paragraph{Equation of state for vertical pressure}
%
\begin{equation}
\tag{5.7.4e}
p
= \left\{
    \begin{tabular}{@{}l@{}}
    some explicit expression which depends on\\
    the explicit assumptions of the model
    \end{tabular}
  \right\}
\end{equation}
%
Possible contributors to vertical pressure are thermal gas pressure, radiation
pressure, magnetic pressure, and turbulent pressure.  Most models to date have
ignored magnetic pressure and turbulent pressure, and have therefore taken
%
\begin{equation}
\tag{5.7.4e$'$}
p = \rho_0\,(T/\mu_{mm}\, m_p)\,T + \tfrac{1}{3}\,b\,T^4,
\end{equation}
%
where $\mu_{mm}$ is the mean molecular weight (0.5 for an ionized hydrogen gas).

\paragraph{Equation for opacity}
%
\begin{equation}
\tag{5.7.4f}
K
= \left\{
    \begin{tabular}{@{}l@{}}
    some explicit expression which depends on\\
    the explicit assumptions of the model
    \end{tabular}
  \right\}
\end{equation}

Of these 6 equations of vertical structure, 3 are differential equations.  They
must be subjected to 3 boundary conditions.  As in the theory of stellar structure,
the boundary conditions at the surface, $z = h$, must be a join to an atmosphere in
which the diffusion approximation for radiative transfer (eq.~5.7.4d$'$) is abandoned.
Alternatively, one can use ``zero-order boundary conditions'' which ignore the
atmosphere---and then ``tack'' an atmosphere onto the model afterwards.  The
zero-order version of the boundary conditions is as follows: (i)~define the surface
of the disk, $z = h$, to be that point at which the density goes to zero
%
\begin{equation}
\tag{5.7.5a}
\rho_0 = 0 \quad \text{at} \quad z = h;
\end{equation}
%
then (ii) temperature must also go to zero at the surface
%
\begin{equation}
\tag{5.7.5b}
T = 0 \quad \text{at} \quad z = h;
\end{equation}
%
(iii)~the flux must vanish on the central plane of the disk
%
\begin{equation}
\tag{5.7.5c}
q^z = 0 \quad \text{at} \quad z = 0;
\end{equation}
%
(iv)~the integrated stress must equal the expression~(5.6.14a) calculated from the
theory of radial structure
%
\begin{equation}
\tag{5.7.5d}
2\int_0^h t_{\varphi r}\,dz = W.
\end{equation}

The equations of \textit{vertical} structure~(5.7.4) and these boundary conditions
automatically guarantee that 3 of the equations of \textit{radial} structure---(5.6.14a,\,b)---
are satisfied.  The vertical structure also provides one with a value for
%
\begin{equation}
\Sigma = 2\int_0^h \rho_0\,dz,
\end{equation}
%
which one can insert into the third radial equation~(5.6.14c) to obtain the mean
radial velocity $\bar{v}^r$.  Then the entire structure, both radial and vertical, is known.


%% ----------------------------------------------------------------
%% 5.8  Approximate Version of Vertical Structure
%% ----------------------------------------------------------------

\subsection{Approximate Version of Vertical Structure}
\label{sec:5.8}

Because of the great current uncertainties about the roles and forms of tur%
bulence and magnetic fields, there is no justification in 1972 for building
sophisticated models of vertical structure.  Therefore, let us solve for the
vertical structure in the same very approximate manner as we used in the
Newtonian treatment of~\S\ref{sec:5.2}.  In particular, let us replace the vertical functions
$\rho_0$, $p$, $t_{\varphi r}$, $T$, and $K$ by their mean values in the disk interior, and let us rewrite
the equations of vertical structure~(5.7.4) in the following vertically-averaged
form.  (i)~\textit{Vertical pressure balance}:
%
\begin{equation}
\tag{5.8.1a}
\bar{h}
= (\bar{p}/\rho_0)^{1/2}\,(r/M)^{1/2}\,2\,\mathscr{Q}^{1/2}\,\mathscr{C}^{-1/2}.
\end{equation}

(ii)~\textit{Source of viscosity}:
%
\begin{equation}
\tag{5.8.1b}
t_{\varphi r} = \alpha\, p
\end{equation}
%
[see eq.~(5.2.32) and preceding discussion].  (iii)~\textit{Energy generation}: we merely
replace $q^z$ by a mean value of $\frac{1}{2}F$, where $F$ is the surface flux as calculated from
the theory of radial structure.  (iv)~\textit{Energy transport}: we assume that radiative
transport dominates over turbulent energy transport, so that
%
\begin{equation}
\tag{5.8.1c}
b\,T^4 = \kappa\,\Sigma\, F.
\end{equation}

(v)~\textit{Equation of state}: Turbulent pressure cannot exceed thermal pressure; if it
did, the turbulence would be supersonic and would quickly dissipate into heat.

% [Page 435]

Magnetic pressure cannot exceed thermal pressure; if it did, the magnetic field
would break free of the disk.  Therefore, we are not far wrong in ignoring
turbulent and magnetic contributions to the pressure, and setting
%
\begin{equation}
\tag{5.8.1d}
p = p^{(\text{rad})} + p^{(\text{gas})},
\end{equation}
%
\begin{equation}
p^{(\text{rad})} = \tfrac{1}{3}\,b\,T^4, \qquad
p^{(\text{gas})} = \rho_0\,(T/\mu_{mm}\, m_p).
\end{equation}

(vi)~\textit{Opacity}: Ignoring line opacity and bound-free opacity, which are less than
or of the order of free-free opacity at the high temperatures of our disk, we set,
%
\begin{equation}
\tag{5.8.1e}
\bar{\kappa} = \kappa_{ff} + \kappa_{es}\,,
\end{equation}
%
\begin{equation}
\bar{\kappa}_{ff}
= (0.64 \times 10^{23})
  \left(\frac{\rho_0}{\text{g/cm}^3}\right)
  \left(\frac{T}{\text{K}}\right)^{\!-7/2}
  \;\frac{\text{cm}^2}{\text{g}}\,,
\qquad
\bar{\kappa}_{es} = 0.40\;\frac{\text{cm}^2}{\text{g}}\,.
\end{equation}


%% ----------------------------------------------------------------
%% 5.9  Explicit Models for Disk
%% ----------------------------------------------------------------

\subsection{Explicit Models for Disk}
\label{sec:5.9}

We shall here combine our approximate equations of vertical structure~(5.8.1)
with our ``exact'' equations of radial structure~(5.6.14) to obtain explicit disk
models.  These models are due originally to Shakura and Sunyaev~(1972)~\cite{ShakuraSunyaev1973}---except
for relativistic corrections, which are due to Thorne~(1973b)~\cite{Thorne1973b}.

In these models we shall express $M$ in units of $3M_\odot$ (a typical black-hole mass);
we shall express $\dot{M}_0$ in units of $10^{17}$\,g\,sec$^{-1}$
(a value that will produce
a total X-ray luminosity typical of the strength of galactic X-ray sources,
$L \sim 10^{37}$\,ergs/sec $\simeq 10^{-9}M_\odot/\text{yr}$);
and we shall express $r$ in units of the radius of the extreme-Kerr
horizon:
%
\begin{equation}
\tag{5.9.1}
M_* \equiv M/10\,M_\odot, \qquad
\dot{M}_{0*} \equiv \dot{M}_0/10^{17}\;\text{g\,sec}^{-1},
\end{equation}
%
\begin{equation}
r_* = r/M = (r/4.4 \times 10^5\;\text{cm})\,M_*^{-1}.
\end{equation}
%
Thus, for galactic X-ray sources, reasonable values are
%
\begin{equation}
\tag{5.9.2a}
M_* \sim \dot{M}_{0*} \sim 1;
\end{equation}
%
for a possible supermassive hole at the center of our Galaxy
($M \sim 3 \times 10^7\,M_\odot$,
$\dot{M}_0 \sim 10^{-4}\,M_\odot/\text{yr}$, see~\S\ref{sec:5.3}), reasonable values are
%
\begin{equation}
\tag{5.9.2b}
M_* \sim 10^7, \qquad \dot{M}_{0*} \sim 10^5.
\end{equation}
%
In addition to the quantities that appear explicitly in the equations of structure,
we shall calculate the optical depth at the center of the disk,
%
\begin{equation}
\tag{5.9.3}
\tau = \kappa\,\Sigma\,;
\end{equation}
%
a rough limit on the strength of the chaotic magnetic field,
%
\begin{equation}
\tag{5.9.4}
B \lesssim (8\pi\, p)^{1/2};
\end{equation}
%
and the characteristic timescale for the gas to move inward from radius~$r$ to the
inner edge of the disk
%
\begin{equation}
\tag{5.9.5}
\Delta(r) = -r/\bar{v}^r.
\end{equation}

The disk can be divided into 3 regions: an ``outer region'' (large radii) in
which gas pressure dominates over radiation pressure, and in which the opacity
is predominantly free-free; a ``middle region'' (smaller radii) in which gas
pressure dominates over radiation pressure, but opacity is predominantly due
to electron scattering; and an ``inner region'' (smallest radii) in which radiation
pressure dominates over gas pressure, and opacity is predominantly due to
electron scattering.  Depending on the size of the mass flux, the inner and middle
regions may or may not exist.  For this reason, we shall use the fully relativistic
equations for the radial structure in all regions, rather than take Newtonian
limits from the beginning for the outer and middle regions.

% [Page 436]

\paragraph{Outer region}

$p = p^{(\text{gas})}$, $\kappa = \kappa_{ff}$.  In this region straightforward algebraic manipulations of
equations~(5.8.1) and (5.6.14) yield the following radial profiles:
%
\begin{align}
\tag{5.9.6}
F &= (0.6 \times 10^{26}\;\text{erg/cm}^2\,\text{sec})\,
    (M_*^{-2}\,\dot{M}_{0*})\,
    r_*^{-3}\,\mathscr{B}^{-1}\,\mathscr{C}^{-1/2}\,\mathscr{Q}, \notag \\[4pt]
%
\Sigma &= (9 \times 10^{5}\;\text{g/cm}^2)\,
    (\alpha^{-4/5}\,M_*^{1/4}\,\dot{M}_{0*}^{7/10})\,
    r_*^{3/4}\,
    \mathscr{B}^{-4/5}\,\mathscr{Q}^{-7/10}\,\mathscr{Q}^{14/5}\,, \notag \\[4pt]
%
h &= (9 \times 10^{2}\;\text{cm})\,
    (\alpha^{-1/10}\,M_*^{3/4}\,\dot{M}_{0*}^{3/20})\,
    r_*^{9/8}\,
    \mathscr{B}^{-3/8}\,\mathscr{Q}^{-11/10}\,\mathscr{C}^{1/2}\,
    \mathscr{Q}^{3/5}\,, \notag \\[4pt]
%
\rho_0 &= (8 \times 10^{-1}\;\text{g/cm}^3)\,
    (\alpha^{-7/10}\,M_*^{-7/10}\,\dot{M}_{0*}^{11/20})\,
    r_*^{-15/8}\notag \\
  &\qquad \times \mathscr{B}^{-2/5}\,\mathscr{Q}^{4/5}\,
    \mathscr{C}^{-1/2}\,\mathscr{Q}^{11/5}\,, \notag \\[4pt]
%
T &= (8 \times 10^{7}\;\text{K})\,
    (\alpha^{-1/5}\,M_*^{-1/5}\,\dot{M}_{0*}^{3/10})\,
    r_*^{-3/4}\,
    \mathscr{B}^{-1/5}\,\mathscr{Q}^{-3/10}\,\mathscr{Q}^{6/5}\,, \notag \\[4pt]
%
\tau_{ff} &= (2 \times 10^{3})\,
    (\alpha^{-4/5}\,M_*^{-7/8}\,\dot{M}_{0*}^{-2/5})\,
    r_*^{-21/16}\,
    \mathscr{B}^{2/5}\,\mathscr{C}^{1/2}\,\mathscr{Q}^{-2/5}\,\mathscr{Q}^{-17/5}\,, \notag \\[4pt]
%
B &\lesssim (7 \times 10^{8}\;\text{G})\,
    (\alpha^{-4/5}\,M_*^{-7/8}\,\dot{M}_{0*}^{2/5})\,
    r_*^{-17/8}\,
    \mathscr{B}^{-2/5}\,\mathscr{C}^{-1/2}\,\mathscr{Q}^{4/5}\,\mathscr{Q}^{-11/5}\,,
\end{align}

% [Page 437]

\begin{align}
\frac{p^{(\text{gas})}}{p^{(\text{rad})}}
&= (0.02)\,
    (\alpha^{-1/10}\,M_*^{-1/10}\,\dot{M}_{0*}^{-1/5})\,
    r_*^{1/2}\,
    \mathscr{Q}^{-4/5}\,\mathscr{C}^{1/2}\,\mathscr{Q}^{-1}\,, \notag \\[4pt]
%
\biggl(\frac{\tau_{ff}}{\tau_{es}}\biggr)
&= (0.6 \times 10^{-5})\,
    (M_*\,\dot{M}_{0*})^{3/2}\,
    \mathscr{B}^{1/2}\,\mathscr{C}^{1/2}\,\mathscr{Q}^{3/2}\,\mathscr{Q}^{-1} \notag \\
&= 50\,(\alpha^{-4/5}\,M_*^{-9/10}\,\dot{M}_{0*}^{1/10})\,
    r_*^{3/2}\,
    \mathscr{B}^{1/4}\,\mathscr{Q}^{-11/20}\,\mathscr{Q}^{1/10}\,, \notag \\[4pt]
%
\Delta(r) &= (2\;\text{sec})\,
    (\alpha^{-4/5}\,M_*^{5/4}\,\dot{M}_{0*}^{-4/5})\,
    r_*^{5/4}\,
    \mathscr{B}^{4/5}\,\mathscr{Q}^{1/10}\,\mathscr{Q}^{-4/5}\,\mathscr{Q}^{-7/10}\,. \notag
\end{align}

It is worth noting that $\mathscr{Q} \to 0$ as $r \to r_{ms}$---i.e., at
the inner edge of the disk the relativistic correction $\mathscr{Q}$ goes to zero smoothly.
$\mathscr{Q} \to 0$, $\mathscr{Q}_{,r} \to 0$ as $r \to r_{ms}$.

The transition to the middle region occurs where $\tau_{ff}/\tau_{es} \sim 1$---i.e., at
%
\begin{equation}
\tag{5.9.7}
r_* = r_{om*}
\equiv 2 \times 10^3\,
(M_*^{-2/3}\,\dot{M}_{0*}^{2/3})\,
\mathscr{B}^{2/3}\,\mathscr{Q}^{-8/15}\,\mathscr{C}^{-1/3}\,\mathscr{Q}^{-1/3}\,\mathscr{Q}^{2/3}\,.
\end{equation}
%
$= 100$ for ``supermassive case.''


\paragraph{Middle region}

$p = p^{(\text{gas})}$, $\kappa = \kappa_{es}$.  In this region the equations of structure~(5.8.1) and~(5.6.14)
yield:
%
\begin{align}
\tag{5.9.8}
F &= (0.6 \times 10^{26}\;\text{erg/cm}^2\,\text{sec})\,
    (M_*^{-2}\,\dot{M}_{0*})\,
    r_*^{-3}\,\mathscr{B}^{-1}\,\mathscr{C}^{-1/2}\,\mathscr{Q}, \notag \\[4pt]
%
\Sigma &= (5 \times 10^{3}\;\text{g/cm}^2)\,
    (\alpha^{-4/5}\,M_*^{2/5}\,\dot{M}_{0*}^{3/5})\,
    r_*^{-3/5}\,
    \mathscr{B}^{-4/5}\,\mathscr{C}^{-1/2}\,\mathscr{Q}^{3/5}\,, \notag \\[4pt]
%
h &= (3 \times 10^{3}\;\text{cm})\,
    (\alpha^{-1/10}\,M_*^{-7/10}\,\dot{M}_{0*}^{1/5})\,
    r_*^{21/20}\,
    \mathscr{B}^{-1/5}\,\mathscr{C}^{1/2}\,\mathscr{Q}^{2/5}\,, \notag \\[4pt]
%
\rho_0 &= (10\;\text{g/cm}^3)\,
    (\alpha^{-7/10}\,M_*^{-7/10}\,\dot{M}_{0*}^{2/5})\,
    r_*^{-33/20}\,
    \mathscr{B}^{-3/5}\,\mathscr{C}^{-1}\,\mathscr{Q}^{1/5}\,, \notag \\[4pt]
%
T &= (3 \times 10^{8}\;\text{K})\,
    (\alpha^{-1/5}\,M_*^{-1/5}\,\dot{M}_{0*}^{2/5})\,
    r_*^{-9/10}\,
    \mathscr{B}^{-2/5}\,\mathscr{Q}^{2/5}\,, \notag \\[4pt]
%
\tau_{es} &= (2 \times 10^{3})\,
    (\alpha^{-4/5}\,M_*^{1/5}\,\dot{M}_{0*}^{3/5})\,
    r_*^{-3/5}\,
    \mathscr{B}^{-4/5}\,\mathscr{C}^{-1/2}\,\mathscr{Q}^{3/5}\,, \notag \\[4pt]
%
B &\lesssim (2 \times 10^{8}\;\text{G})\,
    (\alpha^{-1/5}\,M_*^{-7/10}\,\dot{M}_{0*}^{2/5})\,
    r_*^{-51/40}\,
    \mathscr{B}^{-2/5}\,\mathscr{C}^{-1/2}\,\mathscr{Q}^{1/5}\,, \notag \\[4pt]
%
\frac{p^{(\text{gas})}}{p^{(\text{rad})}}
&= (0.02)\,(\alpha^{-1/10}\,M_*^{-1/10}\,\dot{M}_{0*}^{-1/5})\,
    r_*^{1/2}\,
    \mathscr{Q}^{-4/5}\,\mathscr{C}^{1/2}\,\mathscr{Q}^{-1}\,, \notag \\[4pt]
%
\biggl(\frac{\tau_{ff}}{\tau_{es}}\biggr)
&= (0.6 \times 10^{-5})\,
    (M_*\,\dot{M}_{0*})^{3/2}\,
    \mathscr{B}^{1/2}\,\mathscr{C}^{1/2}\,\mathscr{Q}^{3/2}\,\mathscr{Q}^{-1}\,, \notag \\
&= 50\,(\alpha^{-4/5}\,M_*^{-9/10}\,\dot{M}_{0*}^{1/10})\,
    r_*^{3/2}\,
    \mathscr{B}^{1/4}\,\mathscr{Q}^{-11/20}\,\mathscr{Q}^{1/10}\,, \notag \\[4pt]
%
\Delta(r) &= (0.7\;\text{sec})\,
    (\alpha^{-4/5}\,M_*^{5/8}\,\dot{M}_{0*}^{-3/10})\,
    r_*^{21/16}\,
    \mathscr{B}^{3/5}\,\mathscr{C}^{-1/2}\,\mathscr{Q}^{-3/10}\,. \notag
\end{align}

The transition from the middle region to the inner region occurs where
$p^{(\text{gas})}/p^{(\text{rad})} \sim 1$---i.e., at
%
\begin{equation}
\tag{5.9.9}
r_* = r_{mi*}
\simeq 40\,
(\alpha^{2/21}\,M_*^{2/3}\,\dot{M}_{0*}^{16/21})\,
\mathscr{B}^{-8/21}\,\mathscr{C}^{-2/21}\,\mathscr{Q}^{16/21}
\end{equation}
%
$\simeq 4$ for supermassive case.
